% =====================================================================
% Axona Applications
% David A. Smith, Axona.net
%
% Compile with:  pdflatex Axona-Applications.tex   (run twice for refs)
%            or:  tectonic -X compile Axona-Applications.tex
% Requires: tufte-latex, booktabs, microtype, xcolor (all standard TeX Live)
%
% Typographic style matches the Axona Explainer (tufte-handout).
% =====================================================================

\documentclass[nobib]{tufte-handout}

\usepackage[utf8]{inputenc}
\usepackage{xcolor}
\usepackage{booktabs}
\usepackage{microtype}
\usepackage{fancyhdr}
\usepackage{etoolbox}
\usepackage{needspace}

% Prevent section-heading orphans (heading stranded at the bottom of a
% page).  \needspace reserves vertical room before each heading; if the
% page lacks it, LaTeX breaks before the heading instead of after.
\pretocmd{\section}{\needspace{6\baselineskip}}{}{}
\pretocmd{\subsection}{\needspace{4\baselineskip}}{}{}

\definecolor{rust}{RGB}{185,78,54}
\definecolor{rustsoft}{RGB}{212,168,150}
\definecolor{inkmuted}{RGB}{102,102,102}

% Roman numeral sections to match the Tufte-book typography
\setcounter{secnumdepth}{1}
\renewcommand{\thesection}{\Roman{section}}

\title{Axona Applications}
\author[David A. Smith]{David A. Smith \\[0.3em] \emph{\small Axona.net}}
\date{}

\begin{document}

% fancyhdr-based plain style: bypasses tufte-handout's letterspaced
% running heads (which break on some page-break points) but keeps
% visible page numbers at the foot.
\fancypagestyle{axonaplain}{%
  \fancyhf{}%
  \renewcommand{\headrulewidth}{0pt}%
  \fancyfoot[C]{\thepage}%
}

% --------------------------------------------------------------------
% Custom title page
% --------------------------------------------------------------------
\thispagestyle{empty}
\null\vfill
\begin{center}
{\fontsize{48}{52}\selectfont \textsc{Axona}\par}
\vspace{0.4em}
{\Large\itshape Applications\par}
\vspace{3.5em}
{\large\itshape What runs on Axona today, what gets better if it\\
migrates, and what we are building next.\par}
\vspace{4em}
{\large David A.~Smith\par}
{\itshape\small Axona.net\par}
\vspace{2em}
{\small An Axona Applications Note \quad v0.1 \quad 2026-06-01\par}
\end{center}
\vfill
\null
\newpage

% --------------------------------------------------------------------
% Epigraph page
% --------------------------------------------------------------------
\thispagestyle{empty}
\null\vfill
\begin{flushright}
\itshape The protocol's defaults are the product's features.\par
\end{flushright}
\vfill
\null\newpage

% --------------------------------------------------------------------
% Body opens with page numbers
% --------------------------------------------------------------------
\pagestyle{axonaplain}
\thispagestyle{axonaplain}

\section*{How to Read This}

\newthought{This document has two parts.}

\textbf{Part~I --- Existing Applications: A Deep Dive} examines the
dozen applications that depend on a DHT today, and what changes for
each if its discovery and propagation layer moves to Axona. It is
adapted from the whitepaper chapter of the same name and refreshed to
the \emph{current} state of the protocol: Axona is no longer a
simulator result, it is a live, browser-native network
(\texttt{@axona/protocol} v2.10.0) with an authenticated handshake,
signed and freshness-bounded messages, and a full pub/sub
lifecycle. Where that shipped reality strengthens --- or qualifies ---
the original projection, the text says so.

\textbf{Part~II --- Future Applications: Axona Powered} covers the
applications we are building \emph{on} Axona rather than migrating
\emph{to} it: \texttt{civildefense.io} (the first end-user
application), \textsc{Syzl} (the adaptive social feed), and a
deliberately open pipeline for what comes next.

\marginnote{\emph{On the numbers.} Latency figures in Part~I are
back-of-envelope projections from the 25{,}000-node benchmark in the
whitepaper, applied to each application's typical workload. They are
projections, not field measurements of the migrated application, and
are marked as such throughout. What is \emph{not} a projection is the
protocol surface itself --- the API verbs, the security properties, the
lifecycle semantics --- which is live and testable today.}

% =====================================================================
\section{Existing Applications: A Deep Dive}
% =====================================================================

\newthought{For each application:} what it does, the DHT layer it uses
now, that layer's performance profile, where it falls short of user
expectations, and what changes if it migrates to Axona.\sidenote{%
Streaming-class applications --- those that become possible only with
Axona's latency budget but do not exist today --- are the subject of
Part~II, not this part.}

\subsection{What ``migrate to Axona'' means in 2026}

When the original chapter was written, ``N-DHT / NH-1'' was a
simulator. The claims below now rest on a deployed substrate, which
changes their character from \emph{hypothesis} to \emph{engineering
estimate}:

\begin{itemize}
\item \textbf{Browser-native and live.} Axona was designed for the
WebRTC environment from the start. The reference network runs at
\texttt{axona.net} via the \texttt{axona-peer} browser node and the
\texttt{axona-bridge} signaling broker.\sidenote{The bridge carries
WebRTC offer/answer payloads only and then drops out --- it sees no
application traffic, and any operator can run one. A federated bridge
mesh is on the roadmap.}
\item \textbf{Authenticated, not anonymous-by-accident.} Every peer
link is established through the \texttt{axona/4} authenticated
handshake, with the mesh channel-binding value bound to the DTLS
certificate fingerprints --- so a relaying bridge cannot transparently
intercept ``direct'' peer traffic.
\item \textbf{Signed and fresh.} Every published message is an
Ed25519-signed envelope (format v2) over a domain-tagged core carrying
a per-publisher monotonic sequence number and timestamp, with a
freshness window enforced at ingress. Replay-to-fresh-subscribers is
closed.
\item \textbf{A real lifecycle, not just fan-out.} Beyond
\texttt{pub}/\texttt{sub}, the kernel now exposes \texttt{unsub},
\texttt{pull} (by id \emph{or} latest), \texttt{kill} (creator-only
delete with tombstones), \texttt{unpub} (owner-only topic destroy),
bounded per-topic queues with deterministic eviction, and a hold-time
TTL.\sidenote{24-hour default, 48-hour ceiling, sliding on pull.}
Membership-gated topics are supported via shared-encryption group keys
at the application layer.
\end{itemize}

With that substrate in place, the per-application analysis follows.

\subsection{BitTorrent (Mainline DHT)}

\textbf{What it does.} Mainline DHT is BitTorrent's trackerless
peer-discovery layer: a torrent's \texttt{info\_hash} maps into the
keyspace, swarm peers register under that key, and new peers query the
DHT to find the swarm.

\textbf{Current implementation.} Kademlia, 160-bit IDs, $K{=}8$
buckets, $\alpha{=}3$ parallel queries.\sidenote{The largest production
DHT in the world by node count --- an estimated 10--20 million nodes at
peak.}

\textbf{Performance \& shortfall.} Cold-start \texttt{get\_peers}
latency is typically 5--30~s, dominated by the iterative
$\alpha$-parallel walk through globally-scattered random-ID peers.
Modern clients hide this behind a centralized tracker for the first
discovery round; the DHT is the tracker-resistance \emph{backup}, not
the primary path --- precisely because the cold-start latency is
unacceptable for a primary path.

\textbf{What Axona changes.} Mainline lookups translate directly (same
operation, same data model). Projected lookup latency drops to
$\sim$250--300~ms at 25K nodes --- a 20--100$\times$ improvement that
collapses the user-visible delay below the threshold where tracker
fallback adds value. The qualitative shift: the DHT becomes \emph{the}
discovery layer, not a fallback.

\subsection{IPFS / Filecoin / libp2p}

\textbf{What it does.} Content-addressable storage: content is
referenced by a CID, and the libp2p Kad-DHT routes lookups to peers
that announced they hold it.

\textbf{Current implementation.} libp2p Kad-DHT (Kademlia variant),
default $K{=}20$, $\alpha{=}3$, tens of thousands of nodes.

\textbf{Performance \& shortfall.} CID lookup averages 1--5~s;
provider-record retrieval adds 1--3~s; end-to-end first-byte from cold
cache is 5--15~s --- 25--300$\times$ slower than a centralized HTTP CDN
(50--200~ms). The IPFS gateway tier exists to paper over this with
centralized servers. Separately, \texttt{libp2p-pubsub} (gossipsub) for
IPNS updates carries its own mesh-scaling challenges.

\textbf{What Axona changes.} DHT-layer lookups drop to
$\sim$250--300~ms (projection), making end-to-end first-byte
$\sim$500~ms--1~s --- a 10--20$\times$ improvement that turns
CDN-replacement viability from hypothetical into a real question.
Axona's axonal pub/sub trees would replace gossipsub for IPNS updates
with bounded per-relay fan-out cost. The combination makes IPFS
deployable on browser-class infrastructure for both halves, which it
currently is not.

\subsection{Ethereum (discv5)}

\textbf{What it does.} discv5 is Ethereum's peer-discovery layer --- how
a new node finds peers --- using a Kademlia variant with
verifiable-identity ENRs.

\textbf{Performance \& shortfall.} Cold-start peer acquisition takes
30~s to several minutes (Kademlia walk plus per-peer ENR verification
round-trips). For an on-demand join (a fresh validator, a new searcher)
that is operational friction; mempool gossip via gossipsub adds a
latency-sensitive scaling dimension the static Kademlia table cannot
adapt to.

\textbf{What Axona changes.} Discovery latency drops 10--50$\times$ for
the DHT layer (projection). The locality-aware bootstrap gives new
nodes regionally-sensible initial peer sets --- which matters more for
Ethereum than BitTorrent because mempool propagation is
latency-sensitive. Axona's axonal trees are a natural fit for
``broadcast a transaction to all interested validators in a region,''
addressing the gossipsub limitation directly.\sidenote{Axona's
authenticated handshake maps cleanly onto discv5's identity
requirements --- the migration does not trade away ENR-style verifiable
identity.}

\subsection{Bitcoin Block / Transaction Propagation}

\textbf{What it does.} Bitcoin propagates blocks and transactions
through a peer-to-peer mesh: DNS-seed bootstrap, \texttt{addr}-message
peer discovery, gossip propagation.

\textbf{Performance \& shortfall.} Full-block propagation averages
1--3~s (compact blocks reduce the announce to a few hundred ms plus a
re-request); mempool transaction latency is 50--200~ms. The flooded
gossip is \emph{not} locality-aware --- a US-east $\to$ Asia transaction
may route through Europe --- and MEV concentrates around peers with
privileged low-latency access to miners.

\textbf{What Axona changes.} Locality-aware discovery solves regional
clustering; nodes find nearby peers by default and transactions
propagate along short-path edges. Compact-block propagation becomes
routed pub/sub: every node subscribes to a block-announcement topic,
the miner publishes, the axonal tree fans out at $\sim$100~ms at 50K
nodes (projection). The tree is built once and reused rather than
re-running inventory exchange per peer per round, and peer-to-miner
latency becomes a property of geography rather than peering
relationships --- the privileged-peer problem partially dissolves.

\subsection{Tor Onion Services v3}

\textbf{What it does.} Self-authenticating, location-hidden endpoints;
descriptor lookup is mediated by the HSDir hash ring.

\textbf{Performance \& shortfall.} Rendezvous takes 5--15~s end-to-end,
dominated by Tor circuit setup; the descriptor lookup itself adds
1--3~s. HSDir nodes are infrastructure-tier (a centralization
concession), and descriptor caching is short-lived --- popular services
see fetch storms.

\textbf{What Axona changes.} The descriptor lookup itself drops to
$\sim$300~ms (projection), though the Tor circuit still dominates total
cost. The structural benefit is decentralizing the HSDir tier: Axona's
hop-caching deposits descriptors at any node along the lookup path,
widening the descriptor-fetch tier beyond designated relays. The
directory-authority trust model is unchanged --- Axona widens the cache
tier, it does not replace Tor's anonymity design.

\subsection{Storj / Sia (Decentralized Storage)}

\textbf{What it does.} Files are sharded, encrypted, and replicated
across provider nodes; the DHT mediates shard placement and retrieval.

\textbf{Performance \& shortfall.} Storj shard retrieval is 200--800~ms
for a 4~KB shard (dominated by the connection to a remote provider, not
the lookup); provider discovery at edge is $\sim$1~s. Placement is
\emph{region-blind} --- a Berlin user may hold shards in S\~ao Paulo.

\textbf{What Axona changes.} Locality-aware placement: shards land on
providers in the user's S2 cell with high probability, dropping
retrieval to the regional budget ($\sim$100--200~ms, projection).
Repair coordination (replacing a shard when a provider drops) becomes
an axonal-tree topic with high delivery under churn --- and the new
lifecycle verbs give it explicit semantics: a provider going offline is
a \texttt{kill}/tombstone event, not an inference from silence.

\subsection{Tox / Briar (Peer-to-Peer Messengers)}

\textbf{What it does.} Serverless peer-to-peer messengers. Tox uses a
Kademlia DHT for contact discovery; Briar uses hybrid local + overlay
routing.

\textbf{Performance \& shortfall.} Tox contact-availability checks take
2--10~s; NAT makes initial establishment slower. \emph{Group chat is
expensive} --- every message goes to every member, making large groups
infeasible --- and ``contact came online'' requires polling because
there is no real pub/sub.

\textbf{What Axona changes.} Contact discovery drops to $\sim$250~ms
(projection), comparable to a centralized messenger's offline-status
check. Group chat becomes a pub/sub topic: bounded per-hop fan-out
scaling to thousands of members; ``came online'' becomes an automatic
event. Crucially, the membership and privacy story is now real, not
aspirational: a private group is a \textbf{shared-encryption topic}
(the group key is distributed at the app layer; the protocol carries
opaque ciphertext and gates nothing it shouldn't), and
\texttt{kill}/\texttt{unsub}/\texttt{unpub} give first-class
leave/remove/teardown. This is the first peer-to-peer messenger
architecture that scales to large groups with a centralized service's
latency profile \emph{and} a credible membership model.

\subsection{YaCy (Decentralized Search)}

\textbf{What it does.} A peer-to-peer search engine --- each peer crawls
part of the web; a custom DHT mediates shared indexing and query
distribution.

\textbf{Performance \& shortfall.} Queries take 2--30~s (index lookup
plus query distribution, both $\sim\!\log N$ on a high
$\sim$200--500~ms per-peer constant). Against Google's $\sim$200~ms
median this is unusable interactively --- the UX gap is why YaCy stayed
a research curiosity.

\textbf{What Axona changes.} DHT-layer query distribution drops to
$\sim$250~ms (projection); shard-level result caching via hop caching
makes an interactive YaCy-like engine plausible. Ranked aggregation
across shards is a separate algorithmic problem, and beating Google
needs decades of ranking infrastructure --- but the \emph{latency
floor} stops being the disqualifying factor.

\subsection{Radicle (Decentralized Git Forge)}

\textbf{What it does.} Peer-to-peer git collaboration; repositories
addressed by peer-ID, a DHT layer for repository discovery, custom
routing for pull requests.

\textbf{Performance \& shortfall.} Repository discovery is 1--3~s;
change-set propagation is 5--30~s --- making ``find a project, browse
it'' feel slow versus GitHub and ``watch this repo'' poll-based rather
than push-based.

\textbf{What Axona changes.} Repository discovery drops to
$\sim$300~ms (projection); change-set propagation becomes a
\texttt{commits-on-this-repo} axonal topic with $\sim$200~ms fan-out,
making the forge interactive in the GitHub sense, not just
browsable.\sidenote{The signed-envelope sequence numbers give watchers
a tamper-evident, gap-detectable commit feed for free.}

\subsection{Handshake / ENS (Decentralized DNS)}

\textbf{What it does.} Decentralized name resolution; ownership
recorded on-chain, a DHT-like layer caches recent lookups.

\textbf{Performance \& shortfall.} Hot cached lookups are
$<$100~ms; cold lookups hit the chain and take seconds to minutes, so
cache hit-rate dominates. The cache tier is operated by infrastructure
peers (a centralization tradeoff) because random-Kademlia caching does
not scale.

\textbf{What Axona changes.} Hop caching makes the cache tier
\emph{implicit}: every routed lookup deposits its result at intermediate
nodes, popular names accrue shortcuts naturally, and the explicit
caching infrastructure becomes unnecessary. Cold-cache latency drops
because the routing itself is faster.

\subsection{WebTorrent / WebRTC P2P}

\textbf{What it does.} BitTorrent for browsers over WebRTC data
channels, bridged to Mainline DHT for peer discovery (most browsers
cannot speak Mainline directly).

\textbf{Performance \& shortfall.} Discovery is 5--30~s (Mainline
latency plus bridge overhead). The \emph{Mainline-DHT bridge is a
centralization concession driven entirely by the protocol's inability
to run in browsers} --- the bridge operators are a single point of
failure and surveillance.

\textbf{What Axona changes.} Axona \emph{is} a browser-native DHT ---
the cleanest fit in the catalogue. WebTorrent on Axona eliminates the
\emph{discovery} bridge: peer discovery runs in-browser at the regional
budget (100--300~ms), and the browser becomes a first-class
peer-to-peer participant.\sidenote{Axona still uses a lightweight
\emph{signaling} bridge for WebRTC NAT traversal, but it carries no
application or discovery traffic and drops out once peers are connected
--- a categorically smaller role than the Mainline-DHT bridge it
replaces.} The connection-cap analysis shows this is sustainable on
browser-class infrastructure.

\subsection{Mastodon-Style Federations (No DHT Today)}

\textbf{What it does.} A federation of independent ActivityPub servers;
users follow accounts on other servers, and the home server fetches and
serves federated content.

\textbf{Current implementation.} No DHT. Discovery via WebFinger
(centralized per-domain handle resolution); cross-server delivery via
direct HTTP push from the publisher's home server to each subscriber's
home server.

\textbf{Performance \& shortfall.} Home servers (often small VPS) show
200--2000~ms under load; popular-post federation can lag by minutes due
to per-subscriber-server fan-out. The \emph{home-server-as-bottleneck}
pattern means a viral post on a small instance can take the instance
down, because every other instance's followers fetch from the
publisher's home.

\textbf{What Axona changes.} This is the application that becomes
\emph{architecturally different}. Handle resolution becomes an Axona
lookup (independent of home-server load). Status delivery becomes an
axonal-tree topic per follower-graph cluster: a viral post fans out via
the tree at bounded per-relay cost, never saturating any single home
server. Home servers keep durable storage and WebFinger origin
authority, but the fan-out load moves to the DHT layer --- making small
servers stop being single points of failure for their users' reach.

\subsection{Summary of Improvements}

\begin{table*}[ht]
\centering
\small
\begin{tabular}{@{}p{0.20\linewidth}p{0.36\linewidth}p{0.36\linewidth}@{}}
\toprule
\textbf{Application} & \textbf{Current bottleneck} & \textbf{Axona improvement} \\
\midrule
BitTorrent (Mainline) & 5--30\,s cold lookup & $\sim$300\,ms; primary path possible \\
IPFS / Filecoin & 5--15\,s first-byte & $\sim$1\,s; CDN-replacement viable \\
Ethereum discv5 & Multi-minute join, mempool latency & $\sim$1\,s join, locality-aware mempool \\
Bitcoin propagation & Region-blind, MEV-prone & Locality-aware, axonal-tree announce \\
Tor onion services & DHT minor; HSDir centralization & Wider HSDir tier, faster lookups \\
Storj / Sia & Region-blind shard placement & Locality-aware shards + repair pub/sub \\
Tox / Briar messengers & Slow group chat, polling status & Pub/sub group chat, push status, real membership \\
YaCy & 5--30\,s queries & $\sim$300\,ms DHT layer; interactive plausible \\
Radicle & 5--30\,s changeset propagation & $\sim$200\,ms via axonal pub/sub \\
Handshake / ENS & Cold-cache slow & Implicit caching via hop caching \\
WebTorrent & Mainline-DHT discovery bridge & Discovery bridge eliminated; browser-native \\
Mastodon-style federation & Home-server bottleneck & Decentralized fan-out via axonal trees \\
\bottomrule
\end{tabular}
\caption{Cross-application summary. Latency numbers are projections
from the 25K-node benchmark; production deployments need the
friction-realistic simulator extensions to confirm.}
\label{tab:apps-summary}
\end{table*}

\newthought{The pattern.} The DHT layer is rarely the \emph{only}
bottleneck, but it is consistently \emph{a} bottleneck, and moving to
Axona consistently moves it out of the user-visible critical path. The
biggest beneficiaries are the interactive applications (messengers,
search, dev forges) where DHT latency sits in the user's path; the
least are those where the DHT is already a small fraction of total
latency (Tor circuits, Storj's remote-provider connection). The
application that becomes \emph{architecturally different} is the
federation pattern, because Axona's pub/sub primitive replaces the
home-server fan-out those patterns suffer under.

\textbf{Honest framing.} We keep the federation claim modest: the
architecture is \emph{enabled}, not necessarily preferable in every
dimension. And the latency figures remain projections until measured on
the friction-realistic substrate. What changed since the original
chapter is the floor under those projections --- the protocol is no
longer a simulator artifact but a live, authenticated,
signed-and-fresh network with a real lifecycle.

% =====================================================================
\section{Future Applications: Axona Powered}
% =====================================================================

\newthought{Part~I asked} ``what gets better if an existing app
migrates?'' Part~II asks the opposite: ``what do we build that
\emph{only} Axona makes practical?'' These are applications designed
Axona-native from the first line --- they consume the pub/sub lifecycle,
the geographic locality, the signed-and-fresh envelopes, and the
no-central-operator property as load-bearing features, not
conveniences.

\subsection{civildefense.io --- the flagship}

\textbf{What it is.} A tap-to-report incident map. A user taps a
location to register a concern; reports propagate over anonymous
peer-to-peer with a 24-hour expiry, surfacing on a shared regional map
without any central server holding the data.

\textbf{Why it is Axona-native.} Every primitive the app needs is a
protocol primitive:

\begin{itemize}
\item \textbf{Geographic locality is the product.} The S2 cell prefix
in every nodeId means a report is relevant to, and routes toward, the
people physically near it --- at zero coordination cost. An incident map
is, structurally, a geographically-keyed pub/sub feed, which is exactly
what Axona routes.
\item \textbf{Expiry is a first-class semantic now, not a side
effect.} The original brief leaned on implicit expiry through the
replay-cache LRU. With the v2.10.0 lifecycle, the 24-hour window is the
\textbf{hold-time TTL} --- an explicit, principled bound, so reports age
out deterministically rather than whenever the cache happens to evict
them.
\item \textbf{Signed, fresh, and tamper-evident.} Reports are
Ed25519-signed envelopes with per-publisher sequence numbers and a
freshness window --- so a stale or replayed report cannot be re-injected
at a fresh subscriber, which matters acutely for an incident feed.
\item \textbf{Anonymous but authenticated transport.} The
\texttt{axona/4} handshake plus mesh channel-binding means reports
travel end-to-end without a central operator and without a relaying
bridge being able to read or rewrite them.
\end{itemize}

\textbf{What it still needs from us.} A membership-gated variant
(e.g.\ a verified-responder channel) wants \textbf{Model~3
shared-encryption group keys} --- the group key distributed at the
application layer, the protocol carrying opaque
ciphertext.\sidenote{\texttt{civildefense.io} is one of the two worked
references --- alongside the open-source reference app --- driving that
app-layer group-key design.}

\textbf{Status.} First end-user application; the primitive fit is why
it came together in weeks rather than quarters.

\subsection{\textsc{Syzl} --- the adaptive social feed}

\textbf{What it is.} A swipe-based decentralized social feed where every
gesture trains your network: \emph{swipe right} to \textsc{Syzl} a post
(it reshares to your followers and your connection to the publisher
strengthens), \emph{swipe left} to \textsc{Fyzl} it (it stops, and the
connection weakens). Ranking runs transparently on your own device from
your own swipe history --- there is no algorithm operator and no
server-side ranking.\sidenote{Full product brief:
\texttt{SYZL-Brief.md} in this folder.}

\textbf{Why it is Axona-native.}

\begin{itemize}
\item \textbf{The mechanism mirrors the protocol.} Axona's
neuromorphic routing strengthens useful connections and prunes the
rest; \textsc{Syzl} does the same one layer up, with user attention as
the learning signal. Same algorithm, different layer --- the protocol's
own story retold for end users.
\item \textbf{No central operator, no engagement-maximization
incentive.} No ads, no rage-bait amplification, no shadowbanning;
ranking runs only on-device.
\item \textbf{Verifiable reach without surveillance.} Publishers see
reach (\texttt{publishes} + \texttt{reshare\_count}) via
\texttt{peer.metrics} without learning \emph{who} reshared --- the
privacy property Axona was built for, surfaced at the app layer.
\item \textbf{Connection-sharing as a feature, not a leak.} Because
subscriptions are first-class protocol objects, ``share my graph'' is a
native operation; recipients adopt connections selectively, with
provenance.
\end{itemize}

\textbf{Lifecycle fit.} \textsc{Syzl} is almost entirely
application-layer over \texttt{@axona/protocol}: \texttt{pub}/%
\texttt{sub} for posts, reshares, and connection-list cards;
\texttt{pull} for referenced bodies; \texttt{metrics} for the reach
dashboard; \texttt{unsub} for dropping a rotated-out connection
cleanly. The hill-climb plus simulated-annealing exploration that keeps
the feed alive is $\sim$200 lines on top of \texttt{lookup}/%
\texttt{peers}.

\textbf{\textsc{Syzl} Anywhere.} A companion Manifest-V3 browser
extension that lets a user \textsc{Syzl} any page, selection, or image
in one click, running the user's Axona peer in a background service
worker --- which doubles as keeping the network warm while the user
browses. Roughly six weeks of work, $\sim$70\% code-shared with the PWA
composer.

\subsection{The Pipeline --- others as we invent them}

This section is deliberately open. Part~I is, in effect, a menu:
several of those ``what Axona changes'' analyses describe applications
that are more interesting built \emph{fresh} on Axona than migrated.
Natural near-term candidates, none yet committed:

\begin{itemize}
\item \textbf{A browser-native swarm transport.} WebTorrent's
discovery bridge eliminated --- large-file or live-media distribution
among browser peers with no Mainline-DHT intermediary.
\item \textbf{A large-group peer-to-peer messenger.} Tox/Briar's
group-chat ceiling lifted by pub/sub fan-out plus Model~3
shared-encryption membership.
\item \textbf{A federated-feed fan-out layer.} The Mastodon-style
``viral post on a small instance'' failure removed by moving fan-out to
axonal trees while home servers keep storage and identity.
\end{itemize}

As each of these graduates from idea to product, it gets its own brief
in this folder and a row in Part~II.

\subsection{What the Axona-Powered Applications Share}

Across \texttt{civildefense.io}, \textsc{Syzl}, and the pipeline, the
same four properties keep recurring as the reason the app is
\emph{possible}, not merely \emph{cheaper}:

\begin{enumerate}
\item \textbf{Geographic locality for free} --- the S2 prefix puts
relevance and routing in the same place at zero coordination cost.
\item \textbf{Pub/sub with a real lifecycle} ---
\texttt{pub}/\texttt{sub}/\texttt{unsub}/\texttt{pull}/\texttt{kill}/%
\texttt{unpub}, bounded queues, and hold-time TTL turn ``fan a message
out'' into a managed, expiring, revocable feed.
\item \textbf{Signed, fresh, end-to-end-secure messages} ---
authenticity, freshness, and MITM-resistance are protocol defaults, so
the application never has to bolt them on.
\item \textbf{No central operator} --- ranking, membership, and reach
live on-device and in the mesh, which removes the
surveillance-and-incentive failure mode common to their centralized
counterparts.
\end{enumerate}

These are the same properties Part~I shows existing applications
\emph{reaching toward} through bridges, gateways, and caching tiers. The
forward-looking bet of Part~II is that designing for them from the first
line produces applications the migration path cannot --- because the
protocol's defaults are the product's features.

\end{document}
